\section{Reproducibility and Source Disclosure}
\label{app:code-and-ai-use}

This report is built from the \LaTeX{} sources rooted at
\path{report/final-report.tex}. The solver package is the Julia package
\texttt{StenoticHemodynamics}. The package identity, dependency record, entry
point, and command-line surface are:
\begin{lstlisting}[basicstyle=\ttfamily\footnotesize,breaklines=true]
packages/stenotic-hemodynamics/Project.toml
packages/stenotic-hemodynamics/Manifest.toml
packages/stenotic-hemodynamics/src/StenoticHemodynamics.jl
packages/stenotic-hemodynamics/bin/stenotic-hemodynamics
\end{lstlisting}

\paragraph*{Source snapshot and archive status.}
For a submitted version, the thesis source should be identified by a release tag
or archived source snapshot corresponding to the submitted PDF. This appendix
records the local reproducibility materials available at source-preparation
time; it does not assert that any source copy outside that snapshot is clean,
submission-ready, or identical to a later archived release. The local release
manifest is a source-preparation record:
\begin{lstlisting}[basicstyle=\ttfamily\footnotesize,breaklines=true]
public/reproducibility/release-manifest.json
\end{lstlisting}
This appendix therefore describes local source records without claiming a
public archive. It does not assert a public repository URL, immutable archive
URL, or DOI unless such metadata is added in the release record.

\paragraph*{Archive contents.}
A thesis source snapshot should contain the tracked manuscript source, package
source, dependency manifests or lockfiles, accepted report assets, and
reproducibility records needed to rebuild the submitted copy. It should not
embed local virtual environments, Julia depots, caches, editor state, or other
machine-specific rebuild products such as \path{.venv/}. Those environments are
recreated from the recorded manifests and commands. Full Git history is useful
for a repository-history handoff, for example through a tagged repository or
Git bundle, but it is not required inside an ordinary manuscript source
archive. Large optional raw inputs such as XDMF/HDF5 resolved-flow bundles
should remain separate from the thesis PDF/source archive unless a release plan
explicitly assigns an external data archive, checksum record, and size budget.

\paragraph*{AI-assisted drafting and code work.}
Substantive AI assistance, including OpenAI ChatGPT and Codex, was used during
preparation of this report and repository. That assistance included outlining
and reorganizing prose, drafting and polishing manuscript passages, checking
consistency between the report and local code, suggesting \LaTeX{} and
Julia/Python edits, and helping inspect build, test, and validation output.
AI-generated text or code was treated as draft material, not as an independent
source of evidence or scientific authority. I reviewed, edited, and accept
responsibility for the final manuscript and repository state, including the
mathematical statements, citations, code descriptions, figure and table
descriptions, numerical claims, and conclusions. No AI tool is listed as an
author, credited as a contributor capable of final approval or accountability,
or cited as a primary source.

\paragraph*{Computational environment.}
The repository README states the intended portable execution surface as a macOS
or Linux shell with Julia 1.12 or newer, Python/Pipenv for auxiliary tooling,
and a TeX distribution with \texttt{latexmk} and \texttt{biber}. Concrete
launcher versions, dependency versions, and hardware records belong in the
release manifest or validation logs for the submitted source tree, not as a
standalone manuscript claim.

\paragraph*{Review materials.}
The source inventory, claim register, and review-depth records are maintained as
supplementary files outside the thesis PDF. The bounded claim register for the
main editorial pass and its reader-facing guide are stored at:
\begin{lstlisting}[basicstyle=\ttfamily\footnotesize,breaklines=true]
public/reproducibility/manuscript-claim-register.tsv
public/docs/manuscript-claim-register.md
\end{lstlisting}

\paragraph*{Build command.}
The report build command is:
\begin{lstlisting}[basicstyle=\ttfamily\footnotesize,breaklines=true]
pipenv run ops-build-report --outdir /tmp/masters-report-build --no-sync-final-pdf
\end{lstlisting}
Julia package tests are run through the Python ops validation wrapper:
\begin{lstlisting}[basicstyle=\ttfamily\footnotesize,breaklines=true]
pipenv run ops-julia-check
\end{lstlisting}

\paragraph*{Case-study reconstruction commands.}
The Section~\ref{sec:case-study} verification tables are regenerated with:
\begin{lstlisting}[basicstyle=\ttfamily\footnotesize,breaklines=true]
packages/stenotic-hemodynamics/bin/stenotic-hemodynamics verify mms \
  --nxs 50,100,200,400,800,1600,3200 \
  --dt-values 2e-5,1e-5,5e-6,2.5e-6 \
  --tfinal 0.002 \
  --output-dir report/assets/tables/verification \
  --overwrite

packages/stenotic-hemodynamics/bin/stenotic-hemodynamics verify rest \
  --severities 23,40 \
  --nxs 100,200,400,800,1600,3200 \
  --elapsed-times 0,0.001,0.01,0.1,1.0 \
  --output-dir report/assets/tables/verification \
  --overwrite
\end{lstlisting}
The rest command writes both the drift-history tables and the initial
geometry-rest residual-component CSV/\TeX{} assets consumed in
Section~\ref{subsec:rest-state-drift}.
These timing choices are diagnostic-specific. The MMS run uses the short
manufactured-solution horizon $t=0.002$ s. The zero-forcing rest-state run
intentionally reports elapsed diagnostic times through $t=1.0$ s. The resolved
velocity-comparison workflow below instead targets the imported final-time
sample at $0.9995$ s for cases 77 and 60; this imported-data matching time is
not substituted into the rest-state drift diagnostic.
The additional method-sensitivity rest-state rows in
Table~\ref{tab:rest-method-sensitivity} were generated as scratch experiments
and then rendered into report CSV/\TeX{} fragments:
\begin{lstlisting}[basicstyle=\ttfamily\footnotesize,breaklines=true]
pipenv run ops-experiment --dirty-policy warn \
  --run-id additional-rest-fv-first-order -- \
  verify rest --space fv-first-order \
  --severities 23,40,50 \
  --nxs 400 \
  --elapsed-times 0,0.001,0.01,0.1 \
  --output-dir tmp/simulations/output/additional-evidence/rest-fv-first-order \
  --overwrite

pipenv run ops-experiment --dirty-policy warn \
  --run-id additional-rest-fv-muscl -- \
  verify rest --space fv-muscl \
  --severities 23,40,50 \
  --nxs 400 \
  --elapsed-times 0,0.001,0.01,0.1 \
  --output-dir tmp/simulations/output/additional-evidence/rest-fv-muscl \
  --overwrite

pipenv run ops-experiment --dirty-policy warn \
  --run-id additional-rest-fv-weno3 -- \
  verify rest --space fv-weno3 \
  --severities 23,40,50 \
  --nxs 400 \
  --elapsed-times 0,0.001,0.01,0.1 \
  --output-dir tmp/simulations/output/additional-evidence/rest-fv-weno3 \
  --overwrite

pipenv run ops-experiment --dirty-policy warn \
  --run-id additional-rest-fv-lax-wendroff -- \
  verify rest --space fv-lax-wendroff \
  --severities 23,40,50 \
  --nxs 400 \
  --elapsed-times 0,0.001,0.01,0.1 \
  --output-dir tmp/simulations/output/additional-evidence/rest-fv-lax-wendroff \
  --overwrite

pipenv run ops-render-case-study-evidence-tables
\end{lstlisting}
The Lax--Wendroff rows are retained as a diagnostic-surface capability result:
the current residual decomposition calls the method flux at zero native
timestep, while that flux requires a positive native timestep.
The resolved-velocity comparison uses the 23\% and 40\% stenosis-severity
descriptions in the main text. These correspond to local source IDs 77 and 60,
respectively. The raw 3D fields are optional local inputs rather than archived
thesis assets; the detailed input root and upstream pointer are recorded below.

\paragraph*{Resolved-comparison reconstruction commands.}
When the optional raw 3D inputs are restored locally, the generated comparison
data are regenerated with:
\begin{lstlisting}[basicstyle=\ttfamily\footnotesize,breaklines=true]
packages/stenotic-hemodynamics/bin/stenotic-hemodynamics compare-3d \
  --data-root public/var/data/simulations/canic_case3 \
  --output-dir tmp/simulations/output/3d_comparison/reference_wall_evidence_nx400 \
  --nx 400 \
  --coordinate-mode reference \
  --target-time 0.9995 \
  --time-atol 1e-6 \
  --section-count 200 \
  --profile-slices 1.951,2.451,2.951 \
  --radial-bin-counts 20 \
  --radial-radius-modes current,reference \
  --overwrite \
  --publish-report-assets \
  --report-assets-dir report/assets/data/stenosis-comparison
\end{lstlisting}
The deformed-coordinate companion run changes only the coordinate mode and
scratch output directory:
\begin{lstlisting}[basicstyle=\ttfamily\footnotesize,breaklines=true]
packages/stenotic-hemodynamics/bin/stenotic-hemodynamics compare-3d \
  --data-root public/var/data/simulations/canic_case3 \
  --output-dir tmp/simulations/output/3d_comparison/deformed_wall_evidence_nx400 \
  --nx 400 \
  --coordinate-mode deformed \
  --target-time 0.9995 \
  --time-atol 1e-6 \
  --section-count 200 \
  --profile-slices 1.951,2.451,2.951 \
  --radial-bin-counts 20 \
  --radial-radius-modes current,reference \
  --overwrite \
  --publish-report-assets \
  --report-assets-dir report/assets/data/stenosis-comparison
\end{lstlisting}
These commands regenerate the comparison data in
\path{report/assets/data/stenosis-comparison/}.  The generated file families are:
\begin{lstlisting}[basicstyle=\ttfamily\footnotesize,breaklines=true]
section-quadrature-*
area-audit-*
production-diagnostics-*
node-slab-sensitivity-*
radial-profile-audit-*
\end{lstlisting}
The legacy unsuffixed files in the same directory are compatibility copies of
the reference-coordinate report assets. The compact coordinate-mode and
radial-audit tables are regenerated from those generated data, without raw 3D
inputs, by:
\begin{lstlisting}[basicstyle=\ttfamily\footnotesize,breaklines=true]
pipenv run ops-render-resolved3d-comparison-tables
\end{lstlisting}
The additional timestep-cap sensitivity table in
Table~\ref{tab:t1-time-step-sensitivity} was generated with the live direct
comparison workflow and the same renderer:
\begin{lstlisting}[basicstyle=\ttfamily\footnotesize,breaklines=true]
pipenv run ops-experiment --dirty-policy warn \
  --run-id additional-compare-dt-2e-5 -- \
  compare-3d \
  --data-root public/var/data/simulations/canic_case3 \
  --output-dir tmp/simulations/output/additional-evidence/compare-dt-2e-5 \
  --coordinate-mode reference \
  --nx 400 \
  --dt 2e-5 \
  --target-time 0.9995 \
  --time-atol 1e-6 \
  --section-count 200 \
  --radial-bins 20 \
  --no-svg \
  --overwrite

pipenv run ops-experiment --dirty-policy warn \
  --run-id additional-compare-dt-5e-6 -- \
  compare-3d \
  --data-root public/var/data/simulations/canic_case3 \
  --output-dir tmp/simulations/output/additional-evidence/compare-dt-5e-6 \
  --coordinate-mode reference \
  --nx 400 \
  --dt 5e-6 \
  --target-time 0.9995 \
  --time-atol 1e-6 \
  --section-count 200 \
  --radial-bins 20 \
  --no-svg \
  --overwrite

pipenv run ops-render-case-study-evidence-tables
\end{lstlisting}
A checkout containing these generated comparison data can rebuild the manuscript
and can verify absent raw resolved-3D inputs as an expected skip; it cannot
regenerate the raw-data-dependent comparison assets until the ignored XDMF/HDF5
inputs are restored locally. The remaining timing limitation is transient-history
matching, not the final timestamp.

\paragraph*{Canic Section 4.1 reproducibility comparison.}
The package also provides a separate source-based diagnostic comparison for Canic
et al.\ Section~4.1. It is distinct from the main 23\%/40\% manuscript
comparator above: it reconstructs local 1D outputs, compares section-mean
velocity against restored upstream 3D bundles, emits outlet-gauge pressure
diagnostics, and writes parameter and upstream-provenance comparisons. The 50\%
upstream bundle follows the same per-case final-time rule. Its rows remain
diagnostic comparison output rather than replication evidence unless
parameter, pressure-gauge, and observable-tolerance criteria are also satisfied.
The upstream MATLAB code is treated as external provenance/comparator material
only, not copied implementation code.
\begin{lstlisting}[basicstyle=\ttfamily\footnotesize,breaklines=true]
packages/stenotic-hemodynamics/bin/stenotic-hemodynamics canic-replication section41 \
  --data-root public/var/data/simulations/canic_case3 \
  --output-dir tmp/simulations/output/canic-replication/section41 \
  --coordinate-mode deformed \
  --nx 100 \
  --dt 1e-5 \
  --section-count 200 \
  --radial-sample-count 41 \
  --publish-report-assets \
  --report-assets-dir report/assets \
  --overwrite
\end{lstlisting}
The command emits provenance, parameter-comparison, comparison-summary,
section-comparison, radial-velocity, and compact 3D-diagnostic files. Its
checks intentionally preserve source differences, including the PDF Table~1
Young modulus versus the upstream MATLAB control value and the upstream 50\%
XDMF snapshot time. Table~\ref{tab:canic-section41-summary} summarizes the
generated section-averaged velocity and pressure-diagnostic rows, and
Table~\ref{tab:canic-section41-parameter-audit} lists the source conditions
that bound any Section~4.1 reconstruction or comparison claim. Missing raw
inputs are non-fatal and print:
\begin{lstlisting}[basicstyle=\ttfamily\footnotesize,breaklines=true]
canic_replication_status,skipped_missing_data
\end{lstlisting}
Reviewed Canic derived outputs are stored in
\path{report/assets/data/canic-replication/} and
\path{report/assets/tables/canic-replication/}; the raw XDMF/HDF5 bundles
remain ignored local inputs. The generated Canic fragments used in this revision
were regenerated after the imported-time and outlet-gauge pressure policy
landed, so they are current reproducibility diagnostic assets.

\begin{table}[!htb]
    \centering
    \scriptsize
    \caption{Current Canic Section~4.1 reproducibility comparison boundary.
    Imported times are read from the restored XDMF bundles, and the default
    local target time is the imported time for each case. Explicit global
    final-time overrides are recorded as intentional time mismatches when they
    differ outside the \(10^{-6}\,\mathrm{s}\) tolerance. Pressure diagnostics
    use the common Section~4.1 outlet gauge: the resolved
    \texttt{CrossSectionQuadratureOperator} mean pressure at
    \(z=6\,\mathrm{cm}\) and the corresponding 1D diagnostic outlet pressure
    are subtracted before pressure differences are reported.}
    \label{tab:canic-section41-summary}
    \resizebox{\textwidth}{!}{%
    \begin{tabular}{@{}lllll@{}}
        \toprule
        Case & Imported bundle time & Default local target & Time-alignment role & Pressure role \\
        \midrule
        23\% & \(0.9995\,\mathrm{s}\) & \(0.9995\,\mathrm{s}\) & time-aligned reproducibility comparison & outlet-gauged diagnostic \\
        40\% & \(0.9995\,\mathrm{s}\) & \(0.9995\,\mathrm{s}\) & time-aligned reproducibility comparison & outlet-gauged diagnostic \\
        50\% & \(1.4995\,\mathrm{s}\) & \(1.4995\,\mathrm{s}\) & time-aligned reproducibility comparison & outlet-gauged diagnostic \\
        \bottomrule
    \end{tabular}%
    }
\end{table}

\begin{table}[!htb]
    \centering
    \scriptsize
    \caption{Current Canic Section~4.1 parameter and provenance comparison boundary.
    The rows state the conditions that must accompany any later refreshed
    numerical table.}
    \label{tab:canic-section41-parameter-audit}
    \resizebox{\textwidth}{!}{%
    \begin{tabular}{@{}lll@{}}
        \toprule
        Item & Current policy & Manuscript role \\
        \midrule
        Time alignment & imported, local target, and local completed times must agree within \(10^{-6}\,\mathrm{s}\) & required for replication evidence \\
        Pressure gauge & common Section~4.1 outlet-quadrature gauge is offset-tested and applied & diagnostic pressure differences only \\
        Parameter alignment & PDF Table~1 Young modulus differs from the upstream MATLAB control value & source-model limitation \\
        Raw bundles & local XDMF/HDF5 inputs are restored under \path{public/var/data/simulations/canic_case3/} & optional regeneration input \\
        \bottomrule
    \end{tabular}%
    }
\end{table}

The package also contains the selected quasi-static membrane-FSI comparator. It
solves the radial membrane update on the current lumen geometry and publishes
the report-owned wall profile, fixed-point history, and summary table:
\begin{lstlisting}[basicstyle=\ttfamily\footnotesize,breaklines=true]
packages/stenotic-hemodynamics/bin/stenotic-hemodynamics fsi validate \
  --wall-mode quasi-static \
  --severities 23,40 \
  --meshes 8x2x8,16x4x16 \
  --parallel-workers 0 \
  --output-dir tmp/simulations/output/membrane_fsi_validation \
  --publish-report-assets \
  --report-assets-dir report/assets \
  --overwrite
\end{lstlisting}

\paragraph*{Optional operator and grid-sensitivity evidence.}
The synthetic cross-section operator validation table is regenerated without
raw 3D inputs:
\begin{lstlisting}[basicstyle=\ttfamily\footnotesize,breaklines=true]
packages/stenotic-hemodynamics/bin/stenotic-hemodynamics operator-validation \
  --output-dir report/assets/tables/stenosis-comparison \
  --summary-csv report/assets/data/stenosis-comparison/cross-section-operator-validation.csv \
  --summary-tex report/assets/tables/stenosis-comparison/cross_section_operator_validation.tex \
  --overwrite
\end{lstlisting}
This run uses an in-memory tetrahedron with constant and affine axial fields.
The package tests also include closed-form single-tetrahedron and disjoint
two-tetrahedron cuts with independently computed areas and affine means. These
checks are independent of the optional HDF5/XDMF velocity fields; they do not
validate the available resolved meshes.

When the optional raw 3D inputs are available, the reported grid
sensitivity is:
\begin{lstlisting}[basicstyle=\ttfamily\footnotesize,breaklines=true]
packages/stenotic-hemodynamics/bin/stenotic-hemodynamics compare-3d \
  --data-root public/var/data/simulations/canic_case3 \
  --output-dir tmp/simulations/output/3d_comparison/grid_sensitivity_severity23_severity40 \
  --target-time 0.9995 \
  --time-atol 1e-3 \
  --nxs 200,400,800,1600,3200 \
  --section-count 200 \
  --radial-bins 20 \
  --no-svg \
  --overwrite \
  --grid-summary-csv report/assets/data/stenosis-comparison/grid-sensitivity-summary.csv \
  --grid-summary-tex report/assets/tables/stenosis-comparison/grid_sensitivity_summary.tex
\end{lstlisting}
The manuscript treats this as output-sensitivity evidence for the reported
23\% and 40\% severity section means and physical flows, with $N=1600$ kept as
the intermediate refinement and $N=3200$ as the finest reported grid. It does
not change the matched-data boundary: wall, boundary, material,
current-geometry, and transient-history metadata for the resolved cases remain
unmatched or unpersisted. The target time $0.9995$ s is the imported-bundle
matching time for this workflow; it is distinct from the $t=1.0$ s rest-state
diagnostic and from the one-second native Section~4.1 benchmark target used in
future native resolved-FSI reproduction planning.

Package-benchmark assets are regenerated separately from the PDF:
\begin{lstlisting}[basicstyle=\ttfamily\footnotesize,breaklines=true]
packages/stenotic-hemodynamics/bin/stenotic-hemodynamics benchmark \
  --profile overnight \
  --output-dir tmp/simulations/output/package_benchmark/overnight-YYYYMMDD \
  --overwrite \
  --include-resolved3d \
  --publish-report-assets

pipenv run ops-render-package-benchmark-figures \
  --benchmark-dir tmp/simulations/output/package_benchmark/overnight-YYYYMMDD
\end{lstlisting}
Those rows are secondary numerical diagnostics, not the principal
comparison method.

\paragraph*{Raw 3D inputs and manifest.}
The raw upstream XDMF/HDF5 velocity, pressure, and displacement bundles used for
Section~\ref{sec:resolved-3d-comparison} are not archived in this report
package. The expected local input root is:
\begin{lstlisting}[basicstyle=\ttfamily\footnotesize,breaklines=true]
public/var/data/simulations/canic_case3/
\end{lstlisting}
When available, that directory is restored from this upstream repository,
commit, and subtree:
\begin{lstlisting}[basicstyle=\ttfamily\footnotesize,breaklines=true]
https://github.com/qcutexu/Extended-1D-AQ-system
056a9da2b36b480691f18025d242d2c00f6e7180
case3_all_3d_results/
\end{lstlisting}
For each local case used by the resolved-FSI loaders, the expected companion
files are:
\begin{lstlisting}[basicstyle=\ttfamily\footnotesize,breaklines=true]
velocity.xdmf / velocity.h5
pressure.xdmf / pressure.h5
displace.xdmf / displace.h5
\end{lstlisting}
The report archives generated CSV and PGFPlots-ready static assets, not the raw
fields. The reported comparison data
include the reference-coordinate and deformed-coordinate section summaries
consumed by Section~\ref{subsec:deformed-coordinate-comparison}; wall,
boundary, material, and earlier transient histories remain unresolved unless a
later release entry adds those data.

The release manifest is stored outside the thesis PDF:
\begin{lstlisting}[basicstyle=\ttfamily\footnotesize,breaklines=true]
public/reproducibility/release-manifest.json
\end{lstlisting}
That manifest is release/provenance metadata, not a full report asset
inventory. It records the build command, resolved-input availability, upstream
3D input pointer, comparison reconstruction commands, and the absence of a declared
public repository URL or archival DOI. The report-facing artifact inventory and
command matrix are maintained in \path{public/docs/resolved3d-workflows.md}.
Per-run workflow manifests record SHA-256 output hashes when those workflows are
regenerated.
